\section{Hausdorff's maximal-chain principle}

\import{order/generated-well-order.tex}
\import{order/chain.tex}
\import{set/choice.tex}

Hausdorff's principle is proved from the set-level choice operator by
constructing the canonical well-ordered chain generated by a choice
function.

\begin{theorem}\label{hausdorff_maximal_chain}
    Suppose $Q$ is a binary relation on $P$.
    Suppose $Q$ is transitive and asymmetric.
    Then there exists $C$ such that $C$ is a maximal chain in $P$ with respect to $Q$.
\end{theorem}
\begin{proof}
    Suppose not.
    Let $H=\{C\in\pow{P}\mid \text{$C$ is a chain in $P$ with respect to $Q$}\}$.
    $\emptyset\in\pow{P}$ by \obv{powerset_bottom}.
    $\emptyset$ is a chain in $P$ with respect to $Q$
        by \obv{emptyset_chain}.
    Hence $\emptyset\in H$.
    Thus $H$ is inhabited.

    Show that for all $C\in H$ there exists $p\in P\setminus C$ such that $\cons{p}{C}\in H$.
    \begin{subproof}
        Fix $C\in H$.
        $C$ is a chain in $P$ with respect to $Q$.
        $C$ is not a maximal chain in $P$ with respect to $Q$ by \auto.
        Take $D$ such that $C\subseteq D$ and $D$ is a chain in $P$ with respect to $Q$ and $D\neq C$
            by \auto{maximal_chain_in_relation}.
        $C\subset D$ by \obv{subset}.
        Take $p\in D$ such that $p\notin C$ by \obv{subset_witness}.
        $p\in P\setminus C$ by \auto{chain_in_relation,subseteq,setminus}.
        $\cons{p}{C}\subseteq D$ by \obv{cons_subseteq_intro}.
        Thus $\cons{p}{C}$ is a chain in $P$ with respect to $Q$
            by \obv{subset_of_chain_is_chain}.
        $\cons{p}{C}\subseteq P$ by \obv{chain_in_relation}.
        Hence $\cons{p}{C}\in H$ by \obv{powerset_intro}.
    \end{subproof}

    Let $T=\{w\in H\times P\mid \text{$\snd{w}\notin\fst{w}$ and $\cons{\snd{w}}{\fst{w}}\in H$}\}$.
    Show $T\subseteq H\times P$.
    \begin{subproof}
        It suffices to show that every element of $T$ is an element of $H\times P$
            by \obv{subseteq}.
        Fix $w\in T$.
    \end{subproof}
    $T$ is a relation by \auto{relation,subseteq,times_elem_is_tuple}.
    Show that for all $C\in H$ there exists $p$ such that $C\mathrel{T}p$.
    \begin{subproof}
        Fix $C\in H$.
        Take $p\in P\setminus C$ such that $\cons{p}{C}\in H$.
        $(C,p)\in H\times P$ by \obv{setminus,times_tuple_intro}.
        $\fst{(C,p)}=C$ and $\snd{(C,p)}=p$ by \obv{fst_eq,snd_eq}.
        $\snd{(C,p)}\notin\fst{(C,p)}$ and
            $\cons{\snd{(C,p)}}{\fst{(C,p)}}\in H$
            by \obv{setminus}.
        Thus $(C,p)\in T$.
    \end{subproof}
    Take $f$ such that $f$ is a function on $H$ and for all $C\in H$ we have $C\mathrel{T}f(C)$
        by \obv{axiom_of_choice}.
    Show that for all $C\in H$ we have $f(C)\in P\setminus C$ and $\cons{f(C)}{C}\in H$.
    \begin{subproof}
        Fix $C\in H$.
        $(C,f(C))\in T$.
        $\fst{(C,f(C))}=C$ and $\snd{(C,f(C))}=f(C)$ by \obv{fst_eq,snd_eq}.
        $f(C)\notin C$ and $\cons{f(C)}{C}\in H$.
        $f(C)\in P$ by \obv{subseteq,times_tuple_elim}.
        Follows by \obv{setminus}.
    \end{subproof}

    Let $G=\{p\in H\times\pow{P\times P}\mid \text{$\snd{p}$ is a well ordering generated by $f$ on $\fst{p}$}\}$.
    $(\emptyset,\emptyset)\in G$
        by \auto{times_tuple_intro,powerset_bottom,empty_generated_well_ordering,fst_eq,snd_eq}.
    Thus $G$ is inhabited.
    For all $p\in G$ we have $\snd{p}$ is a well ordering generated by $f$ on $\fst{p}$.

    Let $U=\{\fst{p}\mid p\in G\}$.
    Let $V=\{\snd{p}\mid p\in G\}$.
    Let $B=\unions{U}$.
    Let $S=\unions{V}$.
    $S$ is a well ordering generated by $f$ on $B$ by \obv{union_generated_well_orderings}.

    Show $B$ is a chain in $P$ with respect to $Q$.
    \begin{subproof}
        Show $B\subseteq P$.
        \begin{subproof}
            It suffices to show that every element of $B$ is an element of $P$ by \obv{subseteq}.
            Fix $x\in B$.
            Take $C\in U$ such that $x\in C$ by \obv{unions_iff}.
            Take $p\in G$ such that $C=\fst{p}$.
            $p\in H\times\pow{P\times P}$.
            Thus $C\in H$ by \auto{fst_type}.
            $C\subseteq P$ by \auto{chain_in_relation}.
            Follows by \obv{subseteq}.
        \end{subproof}
        Show that for all $x\in B$ we have for all $y\in B$
            $x=y$ or $x\mathrel{Q}y$ or $y\mathrel{Q}x$.
        \begin{subproof}
            Fix $x\in B$.
            Fix $y\in B$.
            Take $p\in G$ such that $x\in\fst{p}$ by \auto{unions_iff}.
            Take $q\in G$ such that $y\in\fst{q}$ by \auto{unions_iff}.
            \begin{byCase}
                \caseOf{$\fst{p}$ is a common initial segment of $\snd{p}$ on $\fst{p}$ and $\snd{q}$ on $\fst{q}$.}
                    $\fst{p}\subseteq\fst{q}$
                        by \auto{common_initial_segment,inter_elim_right,subseteq}.
                    Thus $x,y\in\fst{q}$ by \obv{subseteq}.
                    $q\in H\times\pow{P\times P}$.
                    Hence $\fst{q}\in H$ by \obv{fst_type}.
                    Thus $\fst{q}$ is a chain in $P$ with respect to $Q$.
                    Follows by \obv{chain_in_relation}.
                \caseOf{$\fst{q}$ is a common initial segment of $\snd{p}$ on $\fst{p}$ and $\snd{q}$ on $\fst{q}$.}
                    $\fst{q}\subseteq\fst{p}$
                        by \auto{common_initial_segment,inter_elim_left,subseteq}.
                    Thus $x,y\in\fst{p}$ by \obv{subseteq}.
                    $p\in H\times\pow{P\times P}$.
                    Hence $\fst{p}\in H$ by \obv{fst_type}.
                    Thus $\fst{p}$ is a chain in $P$ with respect to $Q$.
                    Follows by \obv{chain_in_relation}.
            \end{byCase}
        \end{subproof}
        Follows by \obv{chain_in_relation}.
    \end{subproof}

    $B\subseteq P$ by \obv{chain_in_relation}.
    Hence $B\in\pow{P}$ by \obv{powerset_intro}.
    Thus $B\in H$.
    Let $c=f(B)$.
    $c\in P\setminus B$ and $\cons{c}{B}\in H$ by \auto.
    Let $C=\cons{c}{B}$.
    Let $R=S\union(B\times\{c\})$.
    $(B,c)\in f$ by \auto{function_apply_elim}.
    $R$ is a well ordering generated by $f$ on $C$
        by \obv{generated_well_ordering_adjoin,setminus}.
    $C\in H$.
    Hence $C\in\pow{P}$.
    Thus $C\subseteq P$ by \obv{powerset_elim}.
    $R\subseteq C\times C$ by \auto{generated_well_ordering,strict_well_order_iff}.
    $C\times C\subseteq P\times P$
        by \auto{times_subseteq_left,times_subseteq_right,subseteq_transitive}.
    Hence $R\subseteq P\times P$ by \obv{subseteq_transitive}.
    Thus $R\in\pow{P\times P}$ by \obv{powerset_intro}.
    $(C,R)\in H\times\pow{P\times P}$ by \obv{times_tuple_intro}.
    $\fst{(C,R)}=C$ and $\snd{(C,R)}=R$ by \obv{fst_eq,snd_eq}.
    Hence $(C,R)\in G$ by \auto.
    Thus $C\in U$ by \auto.
    Hence $C\subseteq B$ by \auto{subseteq,unions_iff}.
    Therefore $c\in B$ by \obv{cons_left,subseteq}.
    Contradiction by \obv{setminus_elim_right}.
\end{proof}
